% !TEX root = ../main.tex
% =====================================================================
%  template/chapter_template.tex — MẪU THAM CHIẾU CHO MỌI THÀNH PHẦN
% ---------------------------------------------------------------------
%  * File này KHÔNG nằm trong báo cáo. Dùng để copy cú pháp chuẩn.
%  * Muốn xem bản PDF của mẫu: thêm tạm  % !TEX root = ../main.tex
% =====================================================================
%  template/chapter_template.tex — MẪU THAM CHIẾU CHO MỌI THÀNH PHẦN
% ---------------------------------------------------------------------
%  * File này KHÔNG nằm trong báo cáo. Dùng để copy cú pháp chuẩn.
%  * Muốn xem bản PDF của mẫu: thêm tạm  % !TEX root = ../main.tex
% =====================================================================
%  template/chapter_template.tex — MẪU THAM CHIẾU CHO MỌI THÀNH PHẦN
% ---------------------------------------------------------------------
%  * File này KHÔNG nằm trong báo cáo. Dùng để copy cú pháp chuẩn.
%  * Muốn xem bản PDF của mẫu: thêm tạm  % !TEX root = ../main.tex
% =====================================================================
%  template/chapter_template.tex — MẪU THAM CHIẾU CHO MỌI THÀNH PHẦN
% ---------------------------------------------------------------------
%  * File này KHÔNG nằm trong báo cáo. Dùng để copy cú pháp chuẩn.
%  * Muốn xem bản PDF của mẫu: thêm tạm  \include{template/chapter_template}
%    vào main.tex (KHÔNG commit thay đổi đó).
%  * Thay "cX" bằng tiền tố chương của bạn (c1, c2, ..., c6).
% =====================================================================

% ============================ 1. HEADER ==============================
%  Mỗi file chương PHẢI bắt đầu bằng khối comment header như các file
%  trong chapters/ (Người phụ trách, Tiền tố label, Thư mục ảnh, ...).

\chapter{Mẫu định dạng các thành phần}
\label{chap:cX}

% ======================= 2. CẤU TRÚC MỤC =============================
\section{Tiêu đề mục cấp 1}
\label{sec:cX:muc-cap-1}

\subsection{Tiêu đề mục cấp 2}
\label{sec:cX:muc-cap-2}

\subsubsection{Tiêu đề mục cấp 3 (hạn chế dùng)}

\tieumuc{Tiêu đề nhỏ không đánh số:}
Dùng \verb|\tieumuc{}| cho các ý nhỏ như “Cấu tạo:”, “Nguyên lý hoạt động:”.

% ========================= 3. ĐOẠN VĂN ===============================
\section{Đoạn văn và định dạng chữ}
\label{sec:cX:doan-van}

Mỗi đoạn văn cách nhau bằng MỘT dòng trống. Không dùng \verb|\\| để xuống dòng
cuối đoạn. Dùng \textbf{in đậm} cho thuật ngữ quan trọng khi xuất hiện lần đầu,
\textit{in nghiêng} cho thuật ngữ tiếng Anh, ví dụ: bộ khuếch đại dụng cụ
(\textit{Instrumentation Amplifier}, INA).

% ========================== 4. DANH SÁCH =============================
\section{Danh sách}
\label{sec:cX:danh-sach}

Danh sách không thứ tự:
\begin{itemize}
    \item \textbf{Thân cảm biến:} thanh kim loại đàn hồi.
    \item \textbf{Strain gauge:} điện trở thay đổi theo biến dạng.
\end{itemize}

Danh sách có thứ tự (bước thực hiện):
\begin{enumerate}
    \item Hiệu chỉnh điểm không.
    \item Đặt quả cân chuẩn và ghi nhận kết quả.
\end{enumerate}

Danh sách đánh chữ cái:
\begin{enumerate}[label=\alph*.]
    \item Yêu cầu thứ nhất.
    \item Yêu cầu thứ hai.
\end{enumerate}

% ======================= 5. ĐƠN VỊ (siunitx) =========================
\section{Đại lượng và đơn vị}
\label{sec:cX:don-vi}

Luôn dùng \texttt{\textbackslash qty\{số\}\{đơn vị\}} cho số có đơn vị và \verb|\unit{}| cho đơn vị đứng riêng:
điện áp \qty{3.3}{\volt}, dòng \qty{20}{\milli\ampere}, điện trở \qty{10}{\kilo\ohm},
\qty{3.3}{\ohm}, tụ \qty{100}{\nano\farad}, \qty{1}{\micro\farad}, tần số \qty{1}{\kilo\hertz},
nhiệt độ \qty{25}{\celsius}, khối lượng \qty{2}{\milli\gram}, tỉ lệ \qty{5}{\percent},
độ nhạy \qty{1}{\milli\volt\per\volt}, khoảng \qtyrange{4}{20}{\milli\ampere}.
Đơn vị đứng riêng trong bảng: (\unit{\milli\ampere}).

% ========================= 6. CÔNG THỨC ==============================
\section{Công thức toán}
\label{sec:cX:cong-thuc}

Công thức trong dòng: độ lợi $G = 1 + \frac{\qty{100}{\kilo\ohm}}{R_G}$.

Công thức đứng riêng, CÓ đánh số (khi cần tham chiếu):
\begin{equation}
    G = 1 + \frac{\qty{100}{\kilo\ohm}}{R_G}
    \label{eq:cX:do-loi-ina}
\end{equation}
Tham chiếu: theo \eqref{eq:cX:do-loi-ina}, để $G = 1000$ ta chọn $R_G \approx \qty{100}{\ohm}$.

Chuỗi biến đổi nhiều dòng:
\begin{align}
    \Vout &= I_1 R_f = R_f \frac{\Vin}{R_{\mathrm{in}}} \label{eq:cX:vout} \\
    R_f   &= R_{\mathrm{in}} \cdot RO \cdot \frac{m}{m_{\max}} \label{eq:cX:rf}
\end{align}

Kết quả quan trọng (đóng khung, không đánh số):
\[
    \ketqua{R_f = 1.066 \times m \ (\unit{\ohm})}
\]

Hệ phương trình:
\begin{equation}
    \left\{
    \begin{aligned}
        R_{\mathrm{in}}  &= R_0 + 2R_s \\
        R_{\mathrm{out}} &= R_0 - \frac{R_f^2}{R_0 + 2R_s}
    \end{aligned}
    \right.
    \label{eq:cX:he-pt}
\end{equation}

% =========================== 7. HÌNH ẢNH =============================
\section{Hình ảnh}
\label{sec:cX:hinh-anh}

Một hình (độ rộng mặc định 0.8\verb|\textwidth|), xem Hình~\ref{fig:cX:mot-hinh}:
\hinh[0.6\textwidth]{chuongX/ten_anh.png}{Chú thích ngắn gọn, viết hoa chữ đầu, không dấu chấm cuối}{fig:cX:mot-hinh}

Hai hình cạnh nhau, mỗi hình một số thứ tự:
\haihinh{chuongX/anh_trai.png}{Mô phỏng khi $m = \qty{0}{\kilo\gram}$}{fig:cX:trai}%
        {chuongX/anh_phai.png}{Mô phỏng khi $m = \qty{1}{\kilo\gram}$}{fig:cX:phai}

Nhiều hình con chung một số (subfigure), tham chiếu Hình~\ref{fig:cX:pcb} và Hình~\ref{fig:cX:pcb-top}:
\begin{figure}[H]
    \centering
    \begin{subfigure}[t]{0.47\textwidth}
        \centering
        \chenanh{\textwidth}{chuongX/top.png}
        \caption{Mặt trên}
        \label{fig:cX:pcb-top}
    \end{subfigure}\hfill
    \begin{subfigure}[t]{0.47\textwidth}
        \centering
        \chenanh{\textwidth}{chuongX/bottom.png}
        \caption{Mặt dưới}
        \label{fig:cX:pcb-bottom}
    \end{subfigure}
    \caption{Layout PCB}
    \label{fig:cX:pcb}
\end{figure}

% ============================ 8. BẢNG ================================
\section{Bảng}
\label{sec:cX:bang}

Bảng thông số (2 cột), xem Bảng~\ref{tab:cX:thong-so}. Tên bảng đặt TRÊN bảng:
\begin{table}[H]
    \centering
    \caption{Thông số kỹ thuật của linh kiện}
    \label{tab:cX:thong-so}
    \renewcommand{\arraystretch}{1.3}
    \begin{tabular}{|l|c|}
        \hline
        \textbf{Thông số} & \textbf{Giá trị} \\ \hline
        Điện áp hoạt động & \qtyrange{1.8}{5.5}{\volt} \\ \hline
        Dòng tĩnh         & \qty{50}{\micro\ampere} \\ \hline
    \end{tabular}
\end{table}

Bảng có cột văn bản dài (tự xuống dòng) dùng \texttt{tabularx} với cột \texttt{X}:
\begin{table}[H]
    \centering
    \caption{So sánh các phương án}
    \label{tab:cX:so-sanh}
    \renewcommand{\arraystretch}{1.3}
    \begin{tabularx}{\textwidth}{|c|l|X|}
        \hline
        \textbf{STT} & \textbf{Phương án} & \textbf{Nhận xét} \\ \hline
        1 & Phương án A & Nội dung dài sẽ tự động xuống dòng trong cột X. \\ \hline
        2 & Phương án B & \TODO{} \\ \hline
    \end{tabularx}
\end{table}

Bảng số liệu (tiêu đề cột ghi đơn vị trong ngoặc):
\begin{table}[H]
    \centering
    \caption{Kết quả đo}
    \label{tab:cX:ket-qua}
    \renewcommand{\arraystretch}{1.3}
    \begin{tabular}{|c|c|c|c|}
        \hline
        \textbf{$m$ (\unit{\gram})} & \textbf{$I_{\mathrm{lt}}$ (\unit{\milli\ampere})} &
        \textbf{$I_{\mathrm{tt}}$ (\unit{\milli\ampere})} & \textbf{Sai số (\unit{\percent})} \\ \hline
        0    & 4.00  & 3.96  & 1.00 \\ \hline
        1000 & 20.00 & 19.98 & 0.10 \\ \hline
    \end{tabular}
\end{table}

Bảng gộp ô:
\begin{table}[H]
    \centering
    \caption{Ví dụ gộp ô}
    \label{tab:cX:gop-o}
    \renewcommand{\arraystretch}{1.3}
    \begin{tabular}{|c|c|c|}
        \hline
        \multirow{2}{*}{\textbf{Nhiệt độ}} & \multicolumn{2}{c|}{\textbf{Điện áp ngõ ra (\unit{\volt})}} \\ \cline{2-3}
        & \textbf{Trước bù} & \textbf{Sau bù} \\ \hline
        \qty{0}{\celsius}  & 0.12 & 0.01 \\ \hline
        \qty{50}{\celsius} & 0.15 & 0.02 \\ \hline
    \end{tabular}
\end{table}

% ======================== 9. NHẬN XÉT ===============================
\section{Nhận xét sau kết quả}
\label{sec:cX:nhan-xet}

\nhanxet{Dòng ngõ ra tại \qty{1}{\kilo\gram} đạt \qty{19.98}{\milli\ampere}, sai số
\qty{0.1}{\percent} so với lý thuyết, nằm trong giới hạn $\pm\qty{5}{\percent}$.}

% ========================= 10. TRÍCH DẪN ==============================
\section{Trích dẫn tài liệu}
\label{sec:cX:trich-dan}

Thông số được lấy từ datasheet của nhà sản xuất~\cite{c2_ina333}.
Key trích dẫn phải nằm trong \texttt{references/chuongX.bib} và có tiền tố \texttt{cX\_}.

% ========================== 11. CODE =================================
\section{Mã nguồn}
\label{sec:cX:ma-nguon}

Chỉ dùng chữ KHÔNG DẤU trong code và comment của code.
\begin{lstlisting}[language=C, caption={Doc gia tri ADC}, label={lst:cX:doc-adc}]
uint32_t read_adc(void) {
    // Doc 19 bit tu ADC
    return adc_value & 0x7FFFF;
}
\end{lstlisting}

% ===================== 12. ĐÁNH DẤU CHƯA XONG ========================
\section{Phần chưa hoàn thành}
\label{sec:cX:todo}

\TODO{Ghi chú phần chưa hoàn thành. Xóa hết trước khi merge bản final.}

%    vào main.tex (KHÔNG commit thay đổi đó).
%  * Thay "cX" bằng tiền tố chương của bạn (c1, c2, ..., c6).
% =====================================================================

% ============================ 1. HEADER ==============================
%  Mỗi file chương PHẢI bắt đầu bằng khối comment header như các file
%  trong chapters/ (Người phụ trách, Tiền tố label, Thư mục ảnh, ...).

\chapter{Mẫu định dạng các thành phần}
\label{chap:cX}

% ======================= 2. CẤU TRÚC MỤC =============================
\section{Tiêu đề mục cấp 1}
\label{sec:cX:muc-cap-1}

\subsection{Tiêu đề mục cấp 2}
\label{sec:cX:muc-cap-2}

\subsubsection{Tiêu đề mục cấp 3 (hạn chế dùng)}

\tieumuc{Tiêu đề nhỏ không đánh số:}
Dùng \verb|\tieumuc{}| cho các ý nhỏ như “Cấu tạo:”, “Nguyên lý hoạt động:”.

% ========================= 3. ĐOẠN VĂN ===============================
\section{Đoạn văn và định dạng chữ}
\label{sec:cX:doan-van}

Mỗi đoạn văn cách nhau bằng MỘT dòng trống. Không dùng \verb|\\| để xuống dòng
cuối đoạn. Dùng \textbf{in đậm} cho thuật ngữ quan trọng khi xuất hiện lần đầu,
\textit{in nghiêng} cho thuật ngữ tiếng Anh, ví dụ: bộ khuếch đại dụng cụ
(\textit{Instrumentation Amplifier}, INA).

% ========================== 4. DANH SÁCH =============================
\section{Danh sách}
\label{sec:cX:danh-sach}

Danh sách không thứ tự:
\begin{itemize}
    \item \textbf{Thân cảm biến:} thanh kim loại đàn hồi.
    \item \textbf{Strain gauge:} điện trở thay đổi theo biến dạng.
\end{itemize}

Danh sách có thứ tự (bước thực hiện):
\begin{enumerate}
    \item Hiệu chỉnh điểm không.
    \item Đặt quả cân chuẩn và ghi nhận kết quả.
\end{enumerate}

Danh sách đánh chữ cái:
\begin{enumerate}[label=\alph*.]
    \item Yêu cầu thứ nhất.
    \item Yêu cầu thứ hai.
\end{enumerate}

% ======================= 5. ĐƠN VỊ (siunitx) =========================
\section{Đại lượng và đơn vị}
\label{sec:cX:don-vi}

Luôn dùng \texttt{\textbackslash qty\{số\}\{đơn vị\}} cho số có đơn vị và \verb|\unit{}| cho đơn vị đứng riêng:
điện áp \qty{3.3}{\volt}, dòng \qty{20}{\milli\ampere}, điện trở \qty{10}{\kilo\ohm},
\qty{3.3}{\ohm}, tụ \qty{100}{\nano\farad}, \qty{1}{\micro\farad}, tần số \qty{1}{\kilo\hertz},
nhiệt độ \qty{25}{\celsius}, khối lượng \qty{2}{\milli\gram}, tỉ lệ \qty{5}{\percent},
độ nhạy \qty{1}{\milli\volt\per\volt}, khoảng \qtyrange{4}{20}{\milli\ampere}.
Đơn vị đứng riêng trong bảng: (\unit{\milli\ampere}).

% ========================= 6. CÔNG THỨC ==============================
\section{Công thức toán}
\label{sec:cX:cong-thuc}

Công thức trong dòng: độ lợi $G = 1 + \frac{\qty{100}{\kilo\ohm}}{R_G}$.

Công thức đứng riêng, CÓ đánh số (khi cần tham chiếu):
\begin{equation}
    G = 1 + \frac{\qty{100}{\kilo\ohm}}{R_G}
    \label{eq:cX:do-loi-ina}
\end{equation}
Tham chiếu: theo \eqref{eq:cX:do-loi-ina}, để $G = 1000$ ta chọn $R_G \approx \qty{100}{\ohm}$.

Chuỗi biến đổi nhiều dòng:
\begin{align}
    \Vout &= I_1 R_f = R_f \frac{\Vin}{R_{\mathrm{in}}} \label{eq:cX:vout} \\
    R_f   &= R_{\mathrm{in}} \cdot RO \cdot \frac{m}{m_{\max}} \label{eq:cX:rf}
\end{align}

Kết quả quan trọng (đóng khung, không đánh số):
\[
    \ketqua{R_f = 1.066 \times m \ (\unit{\ohm})}
\]

Hệ phương trình:
\begin{equation}
    \left\{
    \begin{aligned}
        R_{\mathrm{in}}  &= R_0 + 2R_s \\
        R_{\mathrm{out}} &= R_0 - \frac{R_f^2}{R_0 + 2R_s}
    \end{aligned}
    \right.
    \label{eq:cX:he-pt}
\end{equation}

% =========================== 7. HÌNH ẢNH =============================
\section{Hình ảnh}
\label{sec:cX:hinh-anh}

Một hình (độ rộng mặc định 0.8\verb|\textwidth|), xem Hình~\ref{fig:cX:mot-hinh}:
\hinh[0.6\textwidth]{chuongX/ten_anh.png}{Chú thích ngắn gọn, viết hoa chữ đầu, không dấu chấm cuối}{fig:cX:mot-hinh}

Hai hình cạnh nhau, mỗi hình một số thứ tự:
\haihinh{chuongX/anh_trai.png}{Mô phỏng khi $m = \qty{0}{\kilo\gram}$}{fig:cX:trai}%
        {chuongX/anh_phai.png}{Mô phỏng khi $m = \qty{1}{\kilo\gram}$}{fig:cX:phai}

Nhiều hình con chung một số (subfigure), tham chiếu Hình~\ref{fig:cX:pcb} và Hình~\ref{fig:cX:pcb-top}:
\begin{figure}[H]
    \centering
    \begin{subfigure}[t]{0.47\textwidth}
        \centering
        \chenanh{\textwidth}{chuongX/top.png}
        \caption{Mặt trên}
        \label{fig:cX:pcb-top}
    \end{subfigure}\hfill
    \begin{subfigure}[t]{0.47\textwidth}
        \centering
        \chenanh{\textwidth}{chuongX/bottom.png}
        \caption{Mặt dưới}
        \label{fig:cX:pcb-bottom}
    \end{subfigure}
    \caption{Layout PCB}
    \label{fig:cX:pcb}
\end{figure}

% ============================ 8. BẢNG ================================
\section{Bảng}
\label{sec:cX:bang}

Bảng thông số (2 cột), xem Bảng~\ref{tab:cX:thong-so}. Tên bảng đặt TRÊN bảng:
\begin{table}[H]
    \centering
    \caption{Thông số kỹ thuật của linh kiện}
    \label{tab:cX:thong-so}
    \renewcommand{\arraystretch}{1.3}
    \begin{tabular}{|l|c|}
        \hline
        \textbf{Thông số} & \textbf{Giá trị} \\ \hline
        Điện áp hoạt động & \qtyrange{1.8}{5.5}{\volt} \\ \hline
        Dòng tĩnh         & \qty{50}{\micro\ampere} \\ \hline
    \end{tabular}
\end{table}

Bảng có cột văn bản dài (tự xuống dòng) dùng \texttt{tabularx} với cột \texttt{X}:
\begin{table}[H]
    \centering
    \caption{So sánh các phương án}
    \label{tab:cX:so-sanh}
    \renewcommand{\arraystretch}{1.3}
    \begin{tabularx}{\textwidth}{|c|l|X|}
        \hline
        \textbf{STT} & \textbf{Phương án} & \textbf{Nhận xét} \\ \hline
        1 & Phương án A & Nội dung dài sẽ tự động xuống dòng trong cột X. \\ \hline
        2 & Phương án B & \TODO{} \\ \hline
    \end{tabularx}
\end{table}

Bảng số liệu (tiêu đề cột ghi đơn vị trong ngoặc):
\begin{table}[H]
    \centering
    \caption{Kết quả đo}
    \label{tab:cX:ket-qua}
    \renewcommand{\arraystretch}{1.3}
    \begin{tabular}{|c|c|c|c|}
        \hline
        \textbf{$m$ (\unit{\gram})} & \textbf{$I_{\mathrm{lt}}$ (\unit{\milli\ampere})} &
        \textbf{$I_{\mathrm{tt}}$ (\unit{\milli\ampere})} & \textbf{Sai số (\unit{\percent})} \\ \hline
        0    & 4.00  & 3.96  & 1.00 \\ \hline
        1000 & 20.00 & 19.98 & 0.10 \\ \hline
    \end{tabular}
\end{table}

Bảng gộp ô:
\begin{table}[H]
    \centering
    \caption{Ví dụ gộp ô}
    \label{tab:cX:gop-o}
    \renewcommand{\arraystretch}{1.3}
    \begin{tabular}{|c|c|c|}
        \hline
        \multirow{2}{*}{\textbf{Nhiệt độ}} & \multicolumn{2}{c|}{\textbf{Điện áp ngõ ra (\unit{\volt})}} \\ \cline{2-3}
        & \textbf{Trước bù} & \textbf{Sau bù} \\ \hline
        \qty{0}{\celsius}  & 0.12 & 0.01 \\ \hline
        \qty{50}{\celsius} & 0.15 & 0.02 \\ \hline
    \end{tabular}
\end{table}

% ======================== 9. NHẬN XÉT ===============================
\section{Nhận xét sau kết quả}
\label{sec:cX:nhan-xet}

\nhanxet{Dòng ngõ ra tại \qty{1}{\kilo\gram} đạt \qty{19.98}{\milli\ampere}, sai số
\qty{0.1}{\percent} so với lý thuyết, nằm trong giới hạn $\pm\qty{5}{\percent}$.}

% ========================= 10. TRÍCH DẪN ==============================
\section{Trích dẫn tài liệu}
\label{sec:cX:trich-dan}

Thông số được lấy từ datasheet của nhà sản xuất~\cite{c2_ina333}.
Key trích dẫn phải nằm trong \texttt{references/chuongX.bib} và có tiền tố \texttt{cX\_}.

% ========================== 11. CODE =================================
\section{Mã nguồn}
\label{sec:cX:ma-nguon}

Chỉ dùng chữ KHÔNG DẤU trong code và comment của code.
\begin{lstlisting}[language=C, caption={Doc gia tri ADC}, label={lst:cX:doc-adc}]
uint32_t read_adc(void) {
    // Doc 19 bit tu ADC
    return adc_value & 0x7FFFF;
}
\end{lstlisting}

% ===================== 12. ĐÁNH DẤU CHƯA XONG ========================
\section{Phần chưa hoàn thành}
\label{sec:cX:todo}

\TODO{Ghi chú phần chưa hoàn thành. Xóa hết trước khi merge bản final.}

%    vào main.tex (KHÔNG commit thay đổi đó).
%  * Thay "cX" bằng tiền tố chương của bạn (c1, c2, ..., c6).
% =====================================================================

% ============================ 1. HEADER ==============================
%  Mỗi file chương PHẢI bắt đầu bằng khối comment header như các file
%  trong chapters/ (Người phụ trách, Tiền tố label, Thư mục ảnh, ...).

\chapter{Mẫu định dạng các thành phần}
\label{chap:cX}

% ======================= 2. CẤU TRÚC MỤC =============================
\section{Tiêu đề mục cấp 1}
\label{sec:cX:muc-cap-1}

\subsection{Tiêu đề mục cấp 2}
\label{sec:cX:muc-cap-2}

\subsubsection{Tiêu đề mục cấp 3 (hạn chế dùng)}

\tieumuc{Tiêu đề nhỏ không đánh số:}
Dùng \verb|\tieumuc{}| cho các ý nhỏ như “Cấu tạo:”, “Nguyên lý hoạt động:”.

% ========================= 3. ĐOẠN VĂN ===============================
\section{Đoạn văn và định dạng chữ}
\label{sec:cX:doan-van}

Mỗi đoạn văn cách nhau bằng MỘT dòng trống. Không dùng \verb|\\| để xuống dòng
cuối đoạn. Dùng \textbf{in đậm} cho thuật ngữ quan trọng khi xuất hiện lần đầu,
\textit{in nghiêng} cho thuật ngữ tiếng Anh, ví dụ: bộ khuếch đại dụng cụ
(\textit{Instrumentation Amplifier}, INA).

% ========================== 4. DANH SÁCH =============================
\section{Danh sách}
\label{sec:cX:danh-sach}

Danh sách không thứ tự:
\begin{itemize}
    \item \textbf{Thân cảm biến:} thanh kim loại đàn hồi.
    \item \textbf{Strain gauge:} điện trở thay đổi theo biến dạng.
\end{itemize}

Danh sách có thứ tự (bước thực hiện):
\begin{enumerate}
    \item Hiệu chỉnh điểm không.
    \item Đặt quả cân chuẩn và ghi nhận kết quả.
\end{enumerate}

Danh sách đánh chữ cái:
\begin{enumerate}[label=\alph*.]
    \item Yêu cầu thứ nhất.
    \item Yêu cầu thứ hai.
\end{enumerate}

% ======================= 5. ĐƠN VỊ (siunitx) =========================
\section{Đại lượng và đơn vị}
\label{sec:cX:don-vi}

Luôn dùng \texttt{\textbackslash qty\{số\}\{đơn vị\}} cho số có đơn vị và \verb|\unit{}| cho đơn vị đứng riêng:
điện áp \qty{3.3}{\volt}, dòng \qty{20}{\milli\ampere}, điện trở \qty{10}{\kilo\ohm},
\qty{3.3}{\ohm}, tụ \qty{100}{\nano\farad}, \qty{1}{\micro\farad}, tần số \qty{1}{\kilo\hertz},
nhiệt độ \qty{25}{\celsius}, khối lượng \qty{2}{\milli\gram}, tỉ lệ \qty{5}{\percent},
độ nhạy \qty{1}{\milli\volt\per\volt}, khoảng \qtyrange{4}{20}{\milli\ampere}.
Đơn vị đứng riêng trong bảng: (\unit{\milli\ampere}).

% ========================= 6. CÔNG THỨC ==============================
\section{Công thức toán}
\label{sec:cX:cong-thuc}

Công thức trong dòng: độ lợi $G = 1 + \frac{\qty{100}{\kilo\ohm}}{R_G}$.

Công thức đứng riêng, CÓ đánh số (khi cần tham chiếu):
\begin{equation}
    G = 1 + \frac{\qty{100}{\kilo\ohm}}{R_G}
    \label{eq:cX:do-loi-ina}
\end{equation}
Tham chiếu: theo \eqref{eq:cX:do-loi-ina}, để $G = 1000$ ta chọn $R_G \approx \qty{100}{\ohm}$.

Chuỗi biến đổi nhiều dòng:
\begin{align}
    \Vout &= I_1 R_f = R_f \frac{\Vin}{R_{\mathrm{in}}} \label{eq:cX:vout} \\
    R_f   &= R_{\mathrm{in}} \cdot RO \cdot \frac{m}{m_{\max}} \label{eq:cX:rf}
\end{align}

Kết quả quan trọng (đóng khung, không đánh số):
\[
    \ketqua{R_f = 1.066 \times m \ (\unit{\ohm})}
\]

Hệ phương trình:
\begin{equation}
    \left\{
    \begin{aligned}
        R_{\mathrm{in}}  &= R_0 + 2R_s \\
        R_{\mathrm{out}} &= R_0 - \frac{R_f^2}{R_0 + 2R_s}
    \end{aligned}
    \right.
    \label{eq:cX:he-pt}
\end{equation}

% =========================== 7. HÌNH ẢNH =============================
\section{Hình ảnh}
\label{sec:cX:hinh-anh}

Một hình (độ rộng mặc định 0.8\verb|\textwidth|), xem Hình~\ref{fig:cX:mot-hinh}:
\hinh[0.6\textwidth]{chuongX/ten_anh.png}{Chú thích ngắn gọn, viết hoa chữ đầu, không dấu chấm cuối}{fig:cX:mot-hinh}

Hai hình cạnh nhau, mỗi hình một số thứ tự:
\haihinh{chuongX/anh_trai.png}{Mô phỏng khi $m = \qty{0}{\kilo\gram}$}{fig:cX:trai}%
        {chuongX/anh_phai.png}{Mô phỏng khi $m = \qty{1}{\kilo\gram}$}{fig:cX:phai}

Nhiều hình con chung một số (subfigure), tham chiếu Hình~\ref{fig:cX:pcb} và Hình~\ref{fig:cX:pcb-top}:
\begin{figure}[H]
    \centering
    \begin{subfigure}[t]{0.47\textwidth}
        \centering
        \chenanh{\textwidth}{chuongX/top.png}
        \caption{Mặt trên}
        \label{fig:cX:pcb-top}
    \end{subfigure}\hfill
    \begin{subfigure}[t]{0.47\textwidth}
        \centering
        \chenanh{\textwidth}{chuongX/bottom.png}
        \caption{Mặt dưới}
        \label{fig:cX:pcb-bottom}
    \end{subfigure}
    \caption{Layout PCB}
    \label{fig:cX:pcb}
\end{figure}

% ============================ 8. BẢNG ================================
\section{Bảng}
\label{sec:cX:bang}

Bảng thông số (2 cột), xem Bảng~\ref{tab:cX:thong-so}. Tên bảng đặt TRÊN bảng:
\begin{table}[H]
    \centering
    \caption{Thông số kỹ thuật của linh kiện}
    \label{tab:cX:thong-so}
    \renewcommand{\arraystretch}{1.3}
    \begin{tabular}{|l|c|}
        \hline
        \textbf{Thông số} & \textbf{Giá trị} \\ \hline
        Điện áp hoạt động & \qtyrange{1.8}{5.5}{\volt} \\ \hline
        Dòng tĩnh         & \qty{50}{\micro\ampere} \\ \hline
    \end{tabular}
\end{table}

Bảng có cột văn bản dài (tự xuống dòng) dùng \texttt{tabularx} với cột \texttt{X}:
\begin{table}[H]
    \centering
    \caption{So sánh các phương án}
    \label{tab:cX:so-sanh}
    \renewcommand{\arraystretch}{1.3}
    \begin{tabularx}{\textwidth}{|c|l|X|}
        \hline
        \textbf{STT} & \textbf{Phương án} & \textbf{Nhận xét} \\ \hline
        1 & Phương án A & Nội dung dài sẽ tự động xuống dòng trong cột X. \\ \hline
        2 & Phương án B & \TODO{} \\ \hline
    \end{tabularx}
\end{table}

Bảng số liệu (tiêu đề cột ghi đơn vị trong ngoặc):
\begin{table}[H]
    \centering
    \caption{Kết quả đo}
    \label{tab:cX:ket-qua}
    \renewcommand{\arraystretch}{1.3}
    \begin{tabular}{|c|c|c|c|}
        \hline
        \textbf{$m$ (\unit{\gram})} & \textbf{$I_{\mathrm{lt}}$ (\unit{\milli\ampere})} &
        \textbf{$I_{\mathrm{tt}}$ (\unit{\milli\ampere})} & \textbf{Sai số (\unit{\percent})} \\ \hline
        0    & 4.00  & 3.96  & 1.00 \\ \hline
        1000 & 20.00 & 19.98 & 0.10 \\ \hline
    \end{tabular}
\end{table}

Bảng gộp ô:
\begin{table}[H]
    \centering
    \caption{Ví dụ gộp ô}
    \label{tab:cX:gop-o}
    \renewcommand{\arraystretch}{1.3}
    \begin{tabular}{|c|c|c|}
        \hline
        \multirow{2}{*}{\textbf{Nhiệt độ}} & \multicolumn{2}{c|}{\textbf{Điện áp ngõ ra (\unit{\volt})}} \\ \cline{2-3}
        & \textbf{Trước bù} & \textbf{Sau bù} \\ \hline
        \qty{0}{\celsius}  & 0.12 & 0.01 \\ \hline
        \qty{50}{\celsius} & 0.15 & 0.02 \\ \hline
    \end{tabular}
\end{table}

% ======================== 9. NHẬN XÉT ===============================
\section{Nhận xét sau kết quả}
\label{sec:cX:nhan-xet}

\nhanxet{Dòng ngõ ra tại \qty{1}{\kilo\gram} đạt \qty{19.98}{\milli\ampere}, sai số
\qty{0.1}{\percent} so với lý thuyết, nằm trong giới hạn $\pm\qty{5}{\percent}$.}

% ========================= 10. TRÍCH DẪN ==============================
\section{Trích dẫn tài liệu}
\label{sec:cX:trich-dan}

Thông số được lấy từ datasheet của nhà sản xuất~\cite{c2_ina333}.
Key trích dẫn phải nằm trong \texttt{references/chuongX.bib} và có tiền tố \texttt{cX\_}.

% ========================== 11. CODE =================================
\section{Mã nguồn}
\label{sec:cX:ma-nguon}

Chỉ dùng chữ KHÔNG DẤU trong code và comment của code.
\begin{lstlisting}[language=C, caption={Doc gia tri ADC}, label={lst:cX:doc-adc}]
uint32_t read_adc(void) {
    // Doc 19 bit tu ADC
    return adc_value & 0x7FFFF;
}
\end{lstlisting}

% ===================== 12. ĐÁNH DẤU CHƯA XONG ========================
\section{Phần chưa hoàn thành}
\label{sec:cX:todo}

\TODO{Ghi chú phần chưa hoàn thành. Xóa hết trước khi merge bản final.}

%    vào main.tex (KHÔNG commit thay đổi đó).
%  * Thay "cX" bằng tiền tố chương của bạn (c1, c2, ..., c6).
% =====================================================================

% ============================ 1. HEADER ==============================
%  Mỗi file chương PHẢI bắt đầu bằng khối comment header như các file
%  trong chapters/ (Người phụ trách, Tiền tố label, Thư mục ảnh, ...).

\chapter{Mẫu định dạng các thành phần}
\label{chap:cX}

% ======================= 2. CẤU TRÚC MỤC =============================
\section{Tiêu đề mục cấp 1}
\label{sec:cX:muc-cap-1}

\subsection{Tiêu đề mục cấp 2}
\label{sec:cX:muc-cap-2}

\subsubsection{Tiêu đề mục cấp 3 (hạn chế dùng)}

\tieumuc{Tiêu đề nhỏ không đánh số:}
Dùng \verb|\tieumuc{}| cho các ý nhỏ như “Cấu tạo:”, “Nguyên lý hoạt động:”.

% ========================= 3. ĐOẠN VĂN ===============================
\section{Đoạn văn và định dạng chữ}
\label{sec:cX:doan-van}

Mỗi đoạn văn cách nhau bằng MỘT dòng trống. Không dùng \verb|\\| để xuống dòng
cuối đoạn. Dùng \textbf{in đậm} cho thuật ngữ quan trọng khi xuất hiện lần đầu,
\textit{in nghiêng} cho thuật ngữ tiếng Anh, ví dụ: bộ khuếch đại dụng cụ
(\textit{Instrumentation Amplifier}, INA).

% ========================== 4. DANH SÁCH =============================
\section{Danh sách}
\label{sec:cX:danh-sach}

Danh sách không thứ tự:
\begin{itemize}
    \item \textbf{Thân cảm biến:} thanh kim loại đàn hồi.
    \item \textbf{Strain gauge:} điện trở thay đổi theo biến dạng.
\end{itemize}

Danh sách có thứ tự (bước thực hiện):
\begin{enumerate}
    \item Hiệu chỉnh điểm không.
    \item Đặt quả cân chuẩn và ghi nhận kết quả.
\end{enumerate}

Danh sách đánh chữ cái:
\begin{enumerate}[label=\alph*.]
    \item Yêu cầu thứ nhất.
    \item Yêu cầu thứ hai.
\end{enumerate}

% ======================= 5. ĐƠN VỊ (siunitx) =========================
\section{Đại lượng và đơn vị}
\label{sec:cX:don-vi}

Luôn dùng \texttt{\textbackslash qty\{số\}\{đơn vị\}} cho số có đơn vị và \verb|\unit{}| cho đơn vị đứng riêng:
điện áp \qty{3.3}{\volt}, dòng \qty{20}{\milli\ampere}, điện trở \qty{10}{\kilo\ohm},
\qty{3.3}{\ohm}, tụ \qty{100}{\nano\farad}, \qty{1}{\micro\farad}, tần số \qty{1}{\kilo\hertz},
nhiệt độ \qty{25}{\celsius}, khối lượng \qty{2}{\milli\gram}, tỉ lệ \qty{5}{\percent},
độ nhạy \qty{1}{\milli\volt\per\volt}, khoảng \qtyrange{4}{20}{\milli\ampere}.
Đơn vị đứng riêng trong bảng: (\unit{\milli\ampere}).

% ========================= 6. CÔNG THỨC ==============================
\section{Công thức toán}
\label{sec:cX:cong-thuc}

Công thức trong dòng: độ lợi $G = 1 + \frac{\qty{100}{\kilo\ohm}}{R_G}$.

Công thức đứng riêng, CÓ đánh số (khi cần tham chiếu):
\begin{equation}
    G = 1 + \frac{\qty{100}{\kilo\ohm}}{R_G}
    \label{eq:cX:do-loi-ina}
\end{equation}
Tham chiếu: theo \eqref{eq:cX:do-loi-ina}, để $G = 1000$ ta chọn $R_G \approx \qty{100}{\ohm}$.

Chuỗi biến đổi nhiều dòng:
\begin{align}
    \Vout &= I_1 R_f = R_f \frac{\Vin}{R_{\mathrm{in}}} \label{eq:cX:vout} \\
    R_f   &= R_{\mathrm{in}} \cdot RO \cdot \frac{m}{m_{\max}} \label{eq:cX:rf}
\end{align}

Kết quả quan trọng (đóng khung, không đánh số):
\[
    \ketqua{R_f = 1.066 \times m \ (\unit{\ohm})}
\]

Hệ phương trình:
\begin{equation}
    \left\{
    \begin{aligned}
        R_{\mathrm{in}}  &= R_0 + 2R_s \\
        R_{\mathrm{out}} &= R_0 - \frac{R_f^2}{R_0 + 2R_s}
    \end{aligned}
    \right.
    \label{eq:cX:he-pt}
\end{equation}

% =========================== 7. HÌNH ẢNH =============================
\section{Hình ảnh}
\label{sec:cX:hinh-anh}

Một hình (độ rộng mặc định 0.8\verb|\textwidth|), xem Hình~\ref{fig:cX:mot-hinh}:
\hinh[0.6\textwidth]{chuongX/ten_anh.png}{Chú thích ngắn gọn, viết hoa chữ đầu, không dấu chấm cuối}{fig:cX:mot-hinh}

Hai hình cạnh nhau, mỗi hình một số thứ tự:
\haihinh{chuongX/anh_trai.png}{Mô phỏng khi $m = \qty{0}{\kilo\gram}$}{fig:cX:trai}%
        {chuongX/anh_phai.png}{Mô phỏng khi $m = \qty{1}{\kilo\gram}$}{fig:cX:phai}

Nhiều hình con chung một số (subfigure), tham chiếu Hình~\ref{fig:cX:pcb} và Hình~\ref{fig:cX:pcb-top}:
\begin{figure}[H]
    \centering
    \begin{subfigure}[t]{0.47\textwidth}
        \centering
        \chenanh{\textwidth}{chuongX/top.png}
        \caption{Mặt trên}
        \label{fig:cX:pcb-top}
    \end{subfigure}\hfill
    \begin{subfigure}[t]{0.47\textwidth}
        \centering
        \chenanh{\textwidth}{chuongX/bottom.png}
        \caption{Mặt dưới}
        \label{fig:cX:pcb-bottom}
    \end{subfigure}
    \caption{Layout PCB}
    \label{fig:cX:pcb}
\end{figure}

% ============================ 8. BẢNG ================================
\section{Bảng}
\label{sec:cX:bang}

Bảng thông số (2 cột), xem Bảng~\ref{tab:cX:thong-so}. Tên bảng đặt TRÊN bảng:
\begin{table}[H]
    \centering
    \caption{Thông số kỹ thuật của linh kiện}
    \label{tab:cX:thong-so}
    \renewcommand{\arraystretch}{1.3}
    \begin{tabular}{|l|c|}
        \hline
        \textbf{Thông số} & \textbf{Giá trị} \\ \hline
        Điện áp hoạt động & \qtyrange{1.8}{5.5}{\volt} \\ \hline
        Dòng tĩnh         & \qty{50}{\micro\ampere} \\ \hline
    \end{tabular}
\end{table}

Bảng có cột văn bản dài (tự xuống dòng) dùng \texttt{tabularx} với cột \texttt{X}:
\begin{table}[H]
    \centering
    \caption{So sánh các phương án}
    \label{tab:cX:so-sanh}
    \renewcommand{\arraystretch}{1.3}
    \begin{tabularx}{\textwidth}{|c|l|X|}
        \hline
        \textbf{STT} & \textbf{Phương án} & \textbf{Nhận xét} \\ \hline
        1 & Phương án A & Nội dung dài sẽ tự động xuống dòng trong cột X. \\ \hline
        2 & Phương án B & \TODO{} \\ \hline
    \end{tabularx}
\end{table}

Bảng số liệu (tiêu đề cột ghi đơn vị trong ngoặc):
\begin{table}[H]
    \centering
    \caption{Kết quả đo}
    \label{tab:cX:ket-qua}
    \renewcommand{\arraystretch}{1.3}
    \begin{tabular}{|c|c|c|c|}
        \hline
        \textbf{$m$ (\unit{\gram})} & \textbf{$I_{\mathrm{lt}}$ (\unit{\milli\ampere})} &
        \textbf{$I_{\mathrm{tt}}$ (\unit{\milli\ampere})} & \textbf{Sai số (\unit{\percent})} \\ \hline
        0    & 4.00  & 3.96  & 1.00 \\ \hline
        1000 & 20.00 & 19.98 & 0.10 \\ \hline
    \end{tabular}
\end{table}

Bảng gộp ô:
\begin{table}[H]
    \centering
    \caption{Ví dụ gộp ô}
    \label{tab:cX:gop-o}
    \renewcommand{\arraystretch}{1.3}
    \begin{tabular}{|c|c|c|}
        \hline
        \multirow{2}{*}{\textbf{Nhiệt độ}} & \multicolumn{2}{c|}{\textbf{Điện áp ngõ ra (\unit{\volt})}} \\ \cline{2-3}
        & \textbf{Trước bù} & \textbf{Sau bù} \\ \hline
        \qty{0}{\celsius}  & 0.12 & 0.01 \\ \hline
        \qty{50}{\celsius} & 0.15 & 0.02 \\ \hline
    \end{tabular}
\end{table}

% ======================== 9. NHẬN XÉT ===============================
\section{Nhận xét sau kết quả}
\label{sec:cX:nhan-xet}

\nhanxet{Dòng ngõ ra tại \qty{1}{\kilo\gram} đạt \qty{19.98}{\milli\ampere}, sai số
\qty{0.1}{\percent} so với lý thuyết, nằm trong giới hạn $\pm\qty{5}{\percent}$.}

% ========================= 10. TRÍCH DẪN ==============================
\section{Trích dẫn tài liệu}
\label{sec:cX:trich-dan}

Thông số được lấy từ datasheet của nhà sản xuất~\cite{c2_ina333}.
Key trích dẫn phải nằm trong \texttt{references/chuongX.bib} và có tiền tố \texttt{cX\_}.

% ========================== 11. CODE =================================
\section{Mã nguồn}
\label{sec:cX:ma-nguon}

Chỉ dùng chữ KHÔNG DẤU trong code và comment của code.
\begin{lstlisting}[language=C, caption={Doc gia tri ADC}, label={lst:cX:doc-adc}]
uint32_t read_adc(void) {
    // Doc 19 bit tu ADC
    return adc_value & 0x7FFFF;
}
\end{lstlisting}

% ===================== 12. ĐÁNH DẤU CHƯA XONG ========================
\section{Phần chưa hoàn thành}
\label{sec:cX:todo}

\TODO{Ghi chú phần chưa hoàn thành. Xóa hết trước khi merge bản final.}
